\section{Results}\label{sec:results}

This chapter reports the observed results from the frozen Tier 1.5 proof-of-concept. The chapter is limited to factual reporting of the available evidence. Interpretation of what the results mean for the research questions is left to Chapter~\ref{sec:discussion}.

\subsection{Overview of Evaluated System}\label{subsec:results-overview}

The evaluated system consists of the deterministic n8n baseline workflow, the minimal agent-enhanced n8n workflow, the documented safety and validation cases, the CSV-based evaluation harness, and the Raspberry Pi Tier 1.5 validation setup. The local workflow tests used the same two temperature cases defined in \texttt{evaluation/datasets/test-cases.json}: a high-temperature case of \texttt{31.4 C}, expected to produce \texttt{fan\_on}, and a low-temperature case of \texttt{24.5 C}, expected to produce \texttt{fan\_off}.

The local workflow results are recorded in raw CSV files under \texttt{evaluation/results/raw/}. The deterministic run is stored in \texttt{run\_20260429T165733Z\_deterministic.csv}, and the agent-enhanced run is stored in \texttt{run\_20260429T165859Z\_agent.csv}. Processed summaries are stored under \texttt{evaluation/results/processed/}. Raspberry Pi validation evidence is stored under \texttt{evaluation/results/pi-validation/}.

\subsection{Deterministic Baseline Results}\label{subsec:results-deterministic}

The deterministic baseline applies the fixed threshold rule documented in the workflow evidence: temperatures greater than or equal to \texttt{30.0 C} are routed to \texttt{fan\_on}, and lower temperatures are routed to \texttt{fan\_off}. The Step 6 runtime verification notes report that both required branches were verified locally through the n8n test webhook. The final raw deterministic CSV contains one successful row for each defined deterministic test case.

\begin{table}[H]
  \centering
  \caption{Deterministic baseline results from \texttt{evaluation/results/raw/run\_20260429T165733Z\_deterministic.csv}.}
  \label{tab:deterministic-results}
  \begin{tabular}{@{}p{0.28\linewidth}rrrrr@{}}
    \toprule
    Test case & Input (C) & Expected & Observed & HTTP & Latency (ms) \\
    \midrule
    \texttt{deterministic\_high\_temp} & 31.4 & \texttt{fan\_on} & \texttt{fan\_on} & 200 & 480.398 \\
    \texttt{deterministic\_low\_temp} & 24.5 & \texttt{fan\_off} & \texttt{fan\_off} & 200 & 176.778 \\
    \bottomrule
  \end{tabular}
\end{table}

Both deterministic CSV rows have \texttt{success=true}. The observed middleware responses also identify the fan action as simulated. For the high-temperature case, the response contains \texttt{action=fan\_on} and \texttt{fan=on}. For the low-temperature case, the response contains \texttt{action=fan\_off} and \texttt{fan=off}. These rows are the final successful deterministic raw result rows. A previous deterministic production-webhook setup attempt returned HTTP \texttt{404}; that file is stored under \texttt{evaluation/results/raw/dump/} and is treated as setup evidence rather than final action-correctness evidence.

\subsection{Minimal Agent-Enhanced Workflow Results}\label{subsec:results-agent}

The minimal agent-enhanced workflow was evaluated with the same high- and low-temperature inputs. The Step 7 runtime verification evidence reports that the workflow was imported into n8n, connected to a Google Gemini Chat Model, and tested with both defined temperature cases. The evidence includes screenshots for the agent output, parsed action, and resulting middleware fan action. The workflow was configured for stateless execution with no memory node.

For the high-temperature case, the evidence reports that the agent selected \texttt{fan\_on}, produced structured action output containing \texttt{action\_id=fan\_on}, and routed the parsed action to the \texttt{/fan/on} middleware endpoint. For the low-temperature case, the evidence reports that the agent selected \texttt{fan\_off}, produced structured action output containing \texttt{action\_id=fan\_off}, and routed the parsed action to the \texttt{/fan/off} middleware endpoint. The raw agent CSV records the observed middleware actions and response status for these two cases.

\begin{table}[H]
  \centering
  \caption{Minimal agent-enhanced workflow results from \texttt{evaluation/results/raw/run\_20260429T165859Z\_agent.csv}.}
  \label{tab:agent-results}
  \begin{tabular}{@{}p{0.24\linewidth}rrrrr@{}}
    \toprule
    Test case & Input (C) & Expected & Observed & HTTP & Latency (ms) \\
    \midrule
    \texttt{agent\_high\_temp} & 31.4 & \texttt{fan\_on} & \texttt{fan\_on} & 200 & 2668.036 \\
    \texttt{agent\_low\_temp} & 24.5 & \texttt{fan\_off} & \texttt{fan\_off} & 200 & 1995.094 \\
    \bottomrule
  \end{tabular}
\end{table}

Both agent CSV rows have \texttt{success=true}. The CSV file records the final observed middleware action, while the screenshot evidence records the structured agent output and parsed action routing. The result therefore supports only the narrow claim that the minimal agent workflow produced contract-shaped \texttt{fan\_on} and \texttt{fan\_off} actions for the defined cases. It does not provide evidence of broad agent autonomy, model comparison, or multi-agent behavior.

\subsection{Safety and Validation Case Results}\label{subsec:results-contract-safety}

The safety and validation evidence consists of documented cases and expected policy outcomes. It is design/specification evidence, not production-grade runtime enforcement evidence. The safety cases are stored under \texttt{safety\_layer/examples/}. The Step 8 policy defines the allowed actions as \texttt{fan\_on} and \texttt{fan\_off}, the allowed target as \texttt{fan\_1}, and the required fields as \texttt{action\_id}, \texttt{target}, \texttt{reason}, and \texttt{requires\_approval}.

\begin{table}[H]
  \centering
  \caption{Documented safety and validation cases from Step 8.}
  \label{tab:safety-validation-cases}
  \begin{tabular}{@{}p{0.22\linewidth}p{0.26\linewidth}p{0.24\linewidth}p{0.18\linewidth}@{}}
    \toprule
    Case & Input summary & Expected policy outcome & Middleware endpoint \\
    \midrule
    Allowed & Complete \texttt{fan\_on} action for \texttt{fan\_1}; \texttt{requires\_approval=false} & Allow after validation & \texttt{POST /fan/on} \\
    Blocked & \texttt{fan\_on} action missing \texttt{reason} and \texttt{requires\_approval} & Block & None \\
    Risky/approval & \texttt{open\_window} targeting \texttt{unknown\_device}; \texttt{requires\_approval=true} & Reject for direct execution; may be shown for review & None \\
    \bottomrule
  \end{tabular}
\end{table}

The processed file \texttt{evaluation/results/processed/safety\_outcomes.csv} exists, but it currently contains only the header row. It does not contain runtime safety-enforcement rows. For this reason, the safety result in this chapter is reported as documented validation-case evidence only. The repository does not currently prove a full runtime safety framework or production-grade enforcement path.

\subsection{Evaluation Harness Outputs}\label{subsec:results-harness-outputs}

The evaluation harness produced both raw CSV files and processed summary files. The raw CSV files contain one successful deterministic run for each deterministic case and one successful agent run for each agent case. A previous misconfigured deterministic webhook attempt is preserved in \texttt{evaluation/results/raw/dump/} and records HTTP \texttt{404} responses because the production webhook was not registered at the time of that setup attempt.

\begin{table}[H]
  \centering
  \caption{Evaluation harness output files.}
  \label{tab:evaluation-output-files}
  \begin{tabular}{@{}p{0.40\linewidth}p{0.50\linewidth}@{}}
    \toprule
    File or folder & Observed content/status \\
    \midrule
    \texttt{evaluation/results/raw/run\_20260429T165733Z\_deterministic.csv} & Two successful local deterministic workflow rows: high temperature and low temperature. \\
    \texttt{evaluation/results/raw/run\_20260429T165859Z\_agent.csv} & Two successful local agent workflow rows: high temperature and low temperature. \\
    \texttt{evaluation/results/raw/dump/run\_20260429T165702Z\_deterministic.csv} & Setup/misconfigured-webhook evidence with HTTP \texttt{404}; not used as final action-correctness evidence. \\
    \texttt{evaluation/results/processed/summary\_latency.csv} & Processed latency summary exists, but deterministic rows include the earlier setup failure. \\
    \texttt{evaluation/results/processed/summary\_resources.csv} & Processed resource summary exists, with local CPU/RAM fields where available and no local thermal values. \\
    \texttt{evaluation/results/processed/baseline\_vs\_agent.csv} & Baseline-versus-agent summary exists, but deterministic rows include the earlier setup failure. \\
    \texttt{evaluation/results/processed/safety\_outcomes.csv} & Header only; no runtime safety outcome rows. \\
    \bottomrule
  \end{tabular}
\end{table}

The processed summaries are therefore useful as generated artifacts from the aggregation script, but they should be treated carefully. The Step 9 notes state that the earlier deterministic \texttt{404} setup attempt affected the generated deterministic counts. Before final submission, the processed summaries should either be regenerated from the final raw evidence set or explicitly documented as including setup-failure rows. [CHECK VALUE: regenerate \texttt{evaluation/results/processed/*.csv} after excluding setup files, or keep the limitation explicit.]

For the local PC runs, CPU and RAM fields were available in the raw and processed CSV files through best-effort Docker statistics. Local thermal values were recorded as \texttt{not\_available}, which matches the metrics schema expectation for the Windows/PC setup.

[FIGURE NEEDED: result chart/table from CSV]

\subsection{Raspberry Pi Tier 1.5 Validation Results}\label{subsec:results-pi-validation}

The Raspberry Pi validation used a split deployment. n8n stayed on the PC, while the Raspberry Pi ran the Python middleware/action endpoint. The middleware was started on the Raspberry Pi with \texttt{MIDDLEWARE\_HOST=0.0.0.0 python3 -m middleware.api.app}. The PC-hosted n8n production webhook then called the Raspberry Pi middleware over HTTP.

The main full-workflow Pi validation evidence is stored in \texttt{evaluation/results/pi-validation/pi-workflow-run-final.csv}. It records two successful workflow calls: one high-temperature case and one low-temperature case.

\begin{table}[H]
  \centering
  \caption{Raspberry Pi Tier 1.5 full-workflow validation from \texttt{pi-workflow-run-final.csv}.}
  \label{tab:pi-full-workflow-results}
  \begin{tabular}{@{}p{0.28\linewidth}rrrrr@{}}
    \toprule
    Test case & Input (C) & Expected & Observed & HTTP & Latency (ms) \\
    \midrule
    \texttt{deterministic\_high\_temp} & 31.4 & \texttt{fan\_on} & \texttt{fan\_on} & 200 & 437.003 \\
    \texttt{deterministic\_low\_temp} & 24.5 & \texttt{fan\_off} & \texttt{fan\_off} & 200 & 200.206 \\
    \bottomrule
  \end{tabular}
\end{table}

These latency values are full workflow values for the PC-hosted n8n workflow calling the Raspberry Pi middleware. They should be kept separate from direct middleware endpoint latency. The direct endpoint measurements are supporting evidence only because they do not include n8n webhook execution or workflow routing.

\begin{table}[H]
  \centering
  \caption{Supporting Raspberry Pi middleware endpoint and resource observations from \texttt{pi-validation-notes.md} and \texttt{pi-run-01.csv}.}
  \label{tab:pi-supporting-observations}
  \begin{tabular}{@{}p{0.38\linewidth}p{0.32\linewidth}p{0.20\linewidth}@{}}
    \toprule
    Observation & Value & Scope \\
    \midrule
    Direct \texttt{fan\_on} endpoint latency & 135.44 ms & Supporting only \\
    Direct \texttt{fan\_off} endpoint latency & 233.27 ms & Supporting only \\
    Middleware process & \texttt{python3 -m middleware.api.app} & Raspberry Pi middleware \\
    CPU & 0.0 \% & Middleware process observation \\
    MEM & 0.5 \% & Middleware process observation \\
    RSS & 21048 KB & Middleware process observation \\
    Middleware RAM & 20.95 MB & Middleware process observation \\
    Thermal & 49.6 C & Raspberry Pi observation \\
    \bottomrule
  \end{tabular}
\end{table}

The full-workflow CSV retains the Step 9-style fields and records \texttt{thermal\_c=not\_available}. The Raspberry Pi thermal value of \texttt{49.6 C} is documented separately in \texttt{pi-validation-notes.md} and the thermal screenshot evidence. Figure~\ref{fig:results-pi-validation-fan-on-output} shows the selected Raspberry Pi validation output screenshot for the high-temperature fan-on case.

\begin{figure}[H]
  \centering
  \includegraphics[width=\textwidth]{Figures/pi-validation-fan-on-output.png}
  \caption{Raspberry Pi Tier 1.5 validation output for the high-temperature fan-on case.}
  \label{fig:results-pi-validation-fan-on-output}
\end{figure}

The Pi validation result does not show full n8n deployment on the Raspberry Pi. It shows that the PC-hosted n8n workflow could call the Pi-hosted middleware/action endpoint for the defined high- and low-temperature cases. The fan actions remained simulated middleware actions, and no real GPIO hardware was controlled.

\subsection{Summary of Results}\label{subsec:results-summary}

The available results show that the deterministic baseline produced the expected \texttt{fan\_on} and \texttt{fan\_off} actions for the two defined temperature cases in the final raw deterministic CSV. The minimal agent-enhanced workflow produced structured, contract-shaped action output for the same defined cases according to the Step 7 screenshot evidence, and the raw agent CSV records successful \texttt{fan\_on} and \texttt{fan\_off} middleware responses.

The evaluation harness produced raw CSV evidence and processed summary files. The processed summaries are limited because the deterministic processed rows include an earlier setup failure, and the safety outcome summary currently contains only headers. The Raspberry Pi Tier 1.5 validation showed that PC-hosted n8n could call a Raspberry Pi-hosted middleware/action endpoint for the two deterministic temperature cases. These results are limited in scale and should not be treated as broad benchmarking.
